\clearpage
\section{Platt校准：给固定得分再拟合一个概率映射}
\label{sec:diag-platt}
分类器的得分可以很好地排序，却不一定对应事件频率。
二分类Platt校准用两个参数将固定得分$s(\vct x)$映到概率：
\[
 p_{\mathrm{cal}}(\vct x)=\sigma(a\,s(\vct x)+b).
\]
这里采用$+as+b$的Sigmoid约定；也有文献写成$1/(1+e^{As+B})$，相应参数符号相反。
输入是事先定义的得分，不要求它原本就是概率。

\begin{figure}[H]
\centering\begin{tikzpicture}[node distance=7mm]
\node[box,text width=36mm] (train) {训练集\;T\\拟合基础模型};
\node[box,text width=48mm,right=of train] (cal) {校准集\;C\\固定模型产生得分\\拟合映射$a,b$};
\node[box,text width=36mm,right=of cal] (test) {测试集\;E\\固定完整流程\\一次独立评价};
\draw[flow] (train)--(cal);\draw[flow] (cal)--(test);
\end{tikzpicture}

\caption{简化的三个数据角色；交叉拟合可提高数据利用率，但也必须保持每条校准预测的拟合边界。}
\end{figure}
\readfigure{先训练基础模型，再冻结它并在校准集上产生得分，利用校准标签拟合$a,b$。
测试只评价已经固定的完整流程。基础模型若重新拟合，得分分布可能变化，原校准映射也需重新检查。}

沿用前面的8样本例子：两组标签分别为$[0,0,0,1]$与$[0,1,1,1]$，
基础概率为0.1与0.9。若以其logit为得分，则$s=-\log9,+\log9$。
取$a=1/2,b=0$后，$p_{\mathrm{cal}}=0.25,0.75$，正好对应两组观察频率。
这个小算例解释映射，不把校准数据上的改善当作独立测试收益。

\begin{center}
\begin{tabularx}{\textwidth}{P{31mm}Y}
\toprule
方法 & 灵活性与限制\\
\midrule
Sigmoid / Platt & 两参数，数据需求相对较低；真实校准关系若偏离此形状，仍可能不合适。\\
保序回归 Isotonic & 拟合单调的分段映射；更灵活，也更需防止小校准集过拟合与并列概率。\\
\bottomrule
\end{tabularx}
\end{center}

\begin{pitfall}[校准不创造新的排序信息]
$a>0$时该映射严格递增，精确数学下保留排序和ROC-AUC；$a<0$会反转，$a=0$会压成常数。
实际接口是否限制方向、裁剪或产生并列需查实现。概率质量改变不等于AUC必须提高。
少数类太少或部署分布变化时，应重新检查不确定性与校准适用性。
\end{pitfall}

二元形式不能不加说明地当作多分类联合概率；逐类映射后还需明确归一化与评价方式。
校准器和阈值也属于选择流程，应一并纳入验证。
